% Generated draft: requires mathematical and source review.
\documentclass[11pt]{article}
\usepackage[margin=25mm]{geometry}
\usepackage{fontspec}
\setmainfont{Latin Modern Roman}
\usepackage{amsmath,amssymb,mathtools,unicode-math}
\setmathfont{Latin Modern Math}
\usepackage{enumitem,needspace,xurl,hyperref}
\setlength{\parindent}{0pt}
\setlength{\parskip}{5pt}
\setlength{\emergencystretch}{3em}
\newcommand{\N}{\mathbb{N}}
\newcommand{\Z}{\mathbb{Z}}
\newcommand{\Q}{\mathbb{Q}}
\newcommand{\R}{\mathbb{R}}
\newcommand{\C}{\mathbb{C}}
\newcommand{\F}{\mathbb{F}}
\newcommand{\A}{\mathbb{A}}
\newcommand{\PP}{\mathbb{P}}
\newcommand{\E}{\mathbb{E}}
\newcommand{\eps}{\varepsilon}
\newcommand{\dd}{\,\mathrm{d}}
\begin{document}
\section*{UnsolvedMath ID: 30000730}
\textbf{Source problem code:} catalogue entry
\subsection{Definitions and mathematical statement}
Let $p\in\R[z]$ have degree $d>1$ and $p'(0)=\lambda>1$. Let $f$ be the entire solution specified by
\[
 f(\lambda z)=p(f(z)),\qquad f(0)=0,\qquad f'(0)=1.
\]
Evaluation at zero requires $p(0)=0$; this is a compatibility consequence of the supplied equation, not an independently added restriction. Put
\[
 \rho=\log_\lambda d=\frac{\log d}{\log\lambda},\qquad
 N_f(x)=\#\{\xi\in\C:f(\xi)=0,\ |\xi|<x\}\quad(x>0).
\]
The literal counting notation counts distinct complex zeros in the open disk; a multiplicity convention, if intended, must be checked. Let $\mu$ be the source's associated harmonic measure, whose domain and normalization are not supplied.

The conjectural characterization uses the two limits
\[
 \lim_{x\to\infty}x^{-\rho}N_f(x),
 \qquad
 \lim_{t\downarrow0}t^{-\rho}\mu(B(0,t)).
\]
Does existence of either limit characterize polynomials linearly conjugate to a Chebyshev polynomial or a pure power? On the reading that either one separately characterizes the exceptional class, each limit-existence condition is to be equivalent to membership in that class. If the source instead intends only the disjunction of the existence conditions, that is a weaker logical target and requires confirmation. No one-direction implication is substituted for the stated characterization.

The usual polynomial-dynamics meaning of linear conjugacy is affine conjugacy $p=\ell^{-1}\circ q\circ\ell$, with $\ell(z)=az+b$ and $a\ne0$. The source does not state real versus complex conjugacy, whether translations are allowed, the signed Chebyshev convention, or the exact pure-power normal form. These choices matter with the repelling fixed point normalized at zero. Chebyshev polynomials conventionally satisfy $T_d(\cos t)=\cos(dt)$, but no additional sign class is silently included.

The harmonic measure is often that associated with the polynomial's Julia set, viewed from infinity, but the source does not define it. The ball's metric and whether limit existence means a finite value, a finite strictly positive value, or an extended-real value are also unstated. These conventions must be recovered; no zero or infinite limit is silently accepted or excluded.
\subsection{Short English statement}
The conjecture asks whether normalized zero-counting or local harmonic-measure limits occur exactly for Chebyshev-type or pure-power polynomials. It retains both limit formulas, the normalized entire solution, and the unresolved conjugacy, measure, and limit-existence conventions.
\subsection{Sources}
\begin{itemize}
\item Local UnsolvedMath registry, permanent ID 30000730, code catalogue entry, ``Poincar\'e-Function Limits Characterizing Exceptional Polynomials.'' Proposed year 2007; proposer unspecified.
\item \emph{Mini-Workshop: Geometric Measure Theoretic Approaches to Potentials on Fractals and Manifolds} (2008), Oberwolfach Reports, pp.\ 1027--1072, DOI 10.4171/owr/2007/19; \url{https://doi.org/10.4171/owr/2007/19}. The displayed report year 2008 differs from the DOI and proposed year 2007.
\item Supplied ``Poincare-function asymptotics'': \url{https://doi.org/10.1017/S0305004108001564}.
\item Statement status corrected\_verified; the supplied note says official-report content was checked and notation, context, and question form normalized without changing the source claim.
\item Literature triage checked 2026-08-21: open. It says later surveys still state the characterization as conjectural and asymptotic work does not prove the if-and-only-if classification. Exact limit and normalization conventions and current status require independent verification. Created 2026-08-25; updated 2026-09-16.
\item Registry curation by Ulam AI / UnsolvedMath Contributors, CC BY 4.0; underlying report and literature provenance remain as supplied. Unreviewed statement-only draft.
\end{itemize}
\end{document}
